%------------------------------------------------------------------------------%
\documentclass[fleqn,utf8,aspectratio=169,12pt]{beamer}

\usepackage{calculo_varias_variaveis-1}

\title{Coordenadas Polares}

%------------------------------------------------------------------------------%
\begin{document}

\addtitle

%------------------------------------------------------------------------------%
\section{Revisão de Trigonometria}

%------------------------------------------------------------------------------%
\begin{frame}{Relações Trigonométricas no Triângulo Retângulo}
\begin{columns}
\column{0.5\linewidth}
  \begin{align*}
    \onslide<2->{\sin(\theta) & = \frac{\;\;\text{cateto oposto}\;\;}{\text{hipotenusa}} = \frac{\;y\;}{r} \\[5mm]}
    \onslide<3->{\cos(\theta) & = \frac{\text{cateto adjacente}}{\text{hipotenusa}}      = \frac{\;x\;}{r} \\[5mm]}
    \onslide<4->{\tan(\theta) & = \frac{\text{cateto oposto}}{\text{cateto adjacente}}   = \frac{\;y\;}{x}}
  \end{align*}
\column{0.5\linewidth}
  \begin{center}
    \includegraphics{triangulo_retangulo}
  \end{center}
\end{columns}
\end{frame}

%------------------------------------------------------------------------------%
\begin{frame}{Tabela de Seno, Cosseno e Tangente}
  \centering
  \def\arraystretch{2}
  \setlength{\tabcolsep}{2em}
  \begin{tabular}{ccccc}
    \hline
      Graus    &         Radianos         &        $\sin(\theta)$         &        $\cos(\theta)$         &        $\tan(\theta)$         \\ \hline
    $0^\circ$  &            0             &               0               &               1               &               0               \\
    $30^\circ$ & \(\Large\sfrac{\pi}{6}\) &    \(\Large\sfrac{1}{2}\)     & \(\Large\sfrac{\sqrt{3}}{2}\) & \(\Large\sfrac{\sqrt{3}}{3}\) \\
    $45^\circ$ & \(\Large\sfrac{\pi}{4}\) & \(\Large\sfrac{\sqrt{2}}{2}\) & \(\Large\sfrac{\sqrt{2}}{2}\) &               1               \\
    $60^\circ$ & \(\Large\sfrac{\pi}{3}\) & \(\Large\sfrac{\sqrt{3}}{2}\) &    \(\Large\sfrac{1}{2}\)     &         \(\sqrt{3}\)          \\
    $90^\circ$ & \(\Large\sfrac{\pi}{2}\) &               1               &               0               &          $\nexists$           \\ \hline
  \end{tabular}
\end{frame}

%------------------------------------------------------------------------------%
\section{Sistema de Coordenadas Polares}

%------------------------------------------------------------------------------%
\begin{frame}{Coordenadas Polares}
\begin{columns}
\column{0.5\linewidth}
  \begin{itemize}[<+->]
    \item Origem ou polo
    \item Raio inicial
    \item \emph{$\theta$ ângulo orientado} \\
          partindo do raio inicial\\
          (em radianos) \(\theta\in\R\)
    \item \emph{$r$ distância orientada} \\
          partindo da origem, \(r\in\R\)
    \item Coordenadas polares $(r, \theta)$
  \end{itemize}
\column{0.5\linewidth}
\only<1|handout:0>{\includegraphics[page=1]{coordenadas_polares}}%
\only<2|handout:0>{\includegraphics[page=2]{coordenadas_polares}}%
\only<3|handout:0>{\includegraphics[page=3]{coordenadas_polares}}%
\only<4->         {\includegraphics[page=4]{coordenadas_polares}}%
\end{columns}
\end{frame}

%------------------------------------------------------------------------------%
\begin{frame}{Equações em Coordenadas Polares}
\begin{columns}[t]
\column{0.5\linewidth}
  \[
    r = a \qquad a\neq 0
  \]
  \pause
  Círculo centrado na origem e raio $\abs{a}$

  \pause
  \vspace{-0.5\baselineskip}
  \begin{center}
    \includegraphics{equacao_circulo}
  \end{center}
\column{0.5\linewidth}
  \[
    \pause
    \theta=\theta_0
  \]
  \pause
  Reta passando pela origem

  \pause
  \vspace{-0.5\baselineskip}
  \begin{center}
    \includegraphics{equacao_reta}
  \end{center}
\end{columns}
\end{frame}

%------------------------------------------------------------------------------%
\begin{frame}{Relação Entre Coordenadas Polares e Cartesianas}
\begin{columns}
\column{0.2\linewidth}
\begin{align*}
  \onslide<2->{x          & = r\cos(\theta) \\}
  \onslide<3->{y          & = r\sin(\theta) \\[9mm]}
  \onslide<4->{r^2        & = x^2 + y^2     \\}
  \onslide<5->{\tan\theta & = \frac{y}{x}}
\end{align*}
\column{0.7\linewidth}
  \begin{center}
    \onslide<1->{\includegraphics{transformacao}}
  \end{center}
\end{columns}
\end{frame}

%------------------------------------------------------------------------------%
\section{Exemplos}

\include{K-01-Coordenadas_polares-ex1}
\include{K-01-Coordenadas_polares-ex2}
\include{K-01-Coordenadas_polares-ex3}
\include{K-01-Coordenadas_polares-ex4}
\include{K-01-Coordenadas_polares-ex5}
\include{K-01-Coordenadas_polares-ex6}
\include{K-01-Coordenadas_polares-ex7}
\include{K-01-Coordenadas_polares-ex8}
\include{K-01-Coordenadas_polares-ex9}

%------------------------------------------------------------------------------%
\section{Lista Mínima}

%------------------------------------------------------------------------------%
\begin{frame}{Lista Mínima}
  Cálculo Vol. 2 do Thomas 12\textsuperscript{a} ed. -- Seção 11.3
  \begin{enumerate}
    \item Estudar o texto da seção
    \item Resolver os exercícios: 1, 2, 6, 7, 11-15, 27-31, 53-57, 68
  \end{enumerate}

  \vfill
  Atenção: A prova é baseada no livro, não nas apresentações
\end{frame}

%------------------------------------------------------------------------------%
\end{document}
%------------------------------------------------------------------------------%
